\section{Trabalhos relacionados}\label{sec:relacionados}

\subsection{Critérios de seleção}

Esta seção resulta de uma revisão narrativa, e não de uma revisão sistemática.
As buscas foram feitas em novembro de 2026 no Crossref, no OpenAlex, no Google
Scholar e nas páginas das editoras, com combinações dos termos \emph{smart
waste}, \emph{smart bin}, \emph{fill level}, \emph{waste collection routing},
\emph{IoT}, \emph{dynamic routing} e \emph{inventory routing}. Foram
priorizados trabalhos que (i)~usam dados de nível de contêineres para decidir a
coleta e (ii)~avaliam o efeito dessa decisão sobre distância, custo ou nível de
serviço, com preferência para publicações de 2021 a 2026, sem descartar
trabalhos clássicos. Os metadados de todas as referências foram conferidos no
Crossref, no OpenAlex ou na página oficial. Para a maior parte dos trabalhos
comparados, apenas o resumo estava acessível; a Tabela~\ref{tab:relacionados}
registra somente o que consta nesses textos, e células sem informação são
marcadas como ``n.i.''.

\subsection{Panorama}

Os primeiros estudos quantitativos sobre coleta guiada por sensores já
comparavam políticas dinâmicas com rotas fixas. \cite{johansson2006effect}
analisou \num{3300} contêineres de recicláveis com sensores de nível em Malmö,
combinando modelo analítico e simulação de eventos discretos, e relatou
reduções de custo, distância e horas de trabalho com escalonamento e
roteirização dinâmicos. \cite{faccio2011waste} propuseram um modelo
multiobjetivo que usa dados de rastreabilidade em tempo real (sensores
volumétricos, RFID e GPS) e o avaliaram por simulação em uma cidade italiana.
\cite{mes2014inventory} trataram a coleta de contêineres subterrâneos com
sensores como um problema de roteamento de inventário e relataram reduções de
custo de até 40\% com uma heurística antecipatória, cujos parâmetros foram
ajustados por simulação.

Na década seguinte, multiplicaram-se as propostas de roteirização dinâmica. Em
\cite{anagnostopoulos2015robust}, o roteamento considera frota heterogênea e
falhas de caminhões. \cite{akhtar2017backtracking} combinaram CVRP com um
algoritmo de busca por retrocesso e um limiar de nível; o melhor desempenho
ocorreu com limiares de 70\% a 75\%, com redução de \num{36,8}\% na distância.
\cite{hannan2020waste} compararam rotas fixas e variáveis com limiar de 70\% e
relataram ganhos de eficiência e reduções de custo e emissões.
\cite{ramos2018smart} compararam abordagens operacionais para o problema de
coleta inteligente com dados reais. Abordagens preditivas estimam o nível
futuro dos contêineres para antecipar a coleta: \cite{ferrer2019bin} relataram
rotas \num{33,2}\% mais curtas em uma cidade espanhola, e \cite{iosif2026optimizing}
combinaram aprendizado de máquina de curto prazo com VRP dinâmico usando
sensores reais e dados sintéticos.

Trabalhos mais recentes deslocam a atenção para a integração entre dados e
serviço. \cite{barth2023data} usaram 98 contêineres com sensores na Suíça para
estudar a troca entre custo e qualidade de serviço e concluíram que reduções
substanciais de custo são possíveis sem perda de qualidade.
\cite{brouwer2023comparison} compararam formas de monitoramento e mostraram que
a escolha afeta o número de coletas e de transbordamentos. Plataformas como as
de \cite{lozano2018smart} e \cite{cardenasleon2024urban} enfatizam a
infraestrutura de comunicação (LoRaWAN) e a visualização para gestores.
A meta-análise de \cite{maciel2025impact}, com 11 estudos e 21 amostras,
estimou redução média combinada de \num{21,51}\% na distância percorrida e
mostrou que estudos de simulação relatam reduções cerca de três vezes maiores
que as implantações reais.

\begin{table}[ht]
\centering
\caption{Comparação de trabalhos relacionados (informações extraídas dos resumos, salvo indicação). n.i.: não informado.}
\label{tab:relacionados}
\footnotesize
\setlength{\tabcolsep}{3pt}
\renewcommand{\arraystretch}{1.12}
\begin{tabularx}{\textwidth}{@{}>{\raggedright\arraybackslash}p{2.35cm}>{\raggedright\arraybackslash}X>{\raggedright\arraybackslash}p{1.65cm}>{\raggedright\arraybackslash}X>{\raggedright\arraybackslash}X@{}}
\toprule
Trabalho & Sensoriamento e dados & Comunicação & Avaliação e otimização & Limitações para nossa questão \\
\midrule
\cite{johansson2006effect} & nível em \num{3300} contêineres de recicláveis & sem fio (n.i.) & modelo analítico + simulação; políticas dinâmicas vs. rotas fixas & recicláveis; sem falhas de sensor \\
\cite{faccio2011waste} & volume, RFID, GPS & GPRS & VRP multiobjetivo; simulação & ganhos não quantificados no resumo \\
\cite{mes2014inventory} & contêineres subterrâneos com sensores & n.i. & roteamento de inventário; até $-40\%$ de custo & dependência do ajuste de parâmetros \\
\cite{anagnostopoulos2015robust} & dados sintéticos e reais & n.i. & roteamento dinâmico robusto & falhas apenas de caminhão; avaliação qualitativa \\
\cite{akhtar2017backtracking} & nível (limiar) & n.i. & CVRP + busca por retrocesso; $-36{,}8\%$ distância & horizonte de 4 dias; sem teste estatístico declarado \\
\cite{ramos2018smart} & nível em tempo real & n.i. & três abordagens operacionais; caso real & ganho não quantificado no resumo \\
\cite{lozano2018smart} & peso, volume, temperatura & LoRaWAN & rotas dinâmicas; plataforma & contexto rural \\
\cite{ferrer2019bin} & histórico + previsão & n.i. & roteirização a partir de previsão; rotas $-33{,}2\%$ & um caso; comparação com rota da empresa \\
\cite{hannan2020waste} & nível (limiar 70\%) & n.i. & programação linear inteira mista; rotas fixas vs. variáveis & sem incerteza de sensor ou comunicação \\
\cite{barth2023data} & 98 contêineres (campo) & n.i. & troca custo $\times$ serviço & contexto suíço \\
\cite{brouwer2023comparison} & ultrassom vs. observação visual & n.i. & coletas $\times$ transbordamentos & uma empresa \\
\cite{iosif2026optimizing} & sensores reais + sintéticos & n.i. & previsão de curto prazo + VRP dinâmico & sem análise de falhas \\
\cite{maciel2025impact} & --- & --- & meta-análise: $-21{,}5\%$ distância; simulação $-39{,}8\%$ vs. campo $-12{,}4\%$ & só distância; 11 estudos \\
\midrule
Este trabalho & ultrassom simulado com ruído, viés, falhas & MQTT 3.1.1 (broker real e canal emulado) & A, B, C, D pareadas; 30 sementes; fronteiras custo $\times$ serviço; falhas de telemetria & simulação; contêineres hipotéticos; demanda não calibrada em campo \\
\bottomrule
\end{tabularx}
\end{table}

\subsection{Lacuna e posicionamento}

Os trabalhos analisados mostram ganhos consistentes quando a coleta passa a
ser guiada por dados de nível, mas três aspectos aparecem pouco e raramente
juntos nos textos consultados. O primeiro é a modelagem explícita das
imperfeições da telemetria: perda de mensagens, leituras espúrias do
ultrassom, sensores inoperantes e dados desatualizados. \cite{anagnostopoulos2015robust}
tratam falhas de caminhões, não de sensores. O segundo é a comparação
\emph{pareada} de políticas fixa, por limiar, dinâmica e preditiva sob as
mesmas restrições de frota e de jornada, com números aleatórios comuns,
réplicas e testes estatísticos com correção para comparações múltiplas. O
terceiro é a disponibilidade de uma plataforma aberta que inclua a camada de
telemetria e não apenas o otimizador.

Este trabalho se posiciona nessa interseção, de forma incremental. Ele não
propõe um novo algoritmo de roteirização: usa um resolvedor estabelecido e
concentra a contribuição na arquitetura, na modelagem de falhas e no
protocolo experimental. A pergunta implícita na meta-análise de
\cite{maciel2025impact}, isto é, quanto do ganho reportado em simulação resiste a
condições menos favoráveis, orientou o desenho dos experimentos E2 e E3. Como a
busca não foi sistemática e vários trabalhos foram lidos apenas pelo resumo,
não se afirma que essa combinação seja inédita.
